% concatenated from non_compact_f4.tex

\documentclass[11pt]{amsart}
\usepackage{newlfont,amsmath,mathrsfs,enumerate,amssymb,array,amscd,xcolor,nopageno}

\usepackage{graphicx}
%\usepackage{amsmath,amssymb,enumerate}
\usepackage[all]{xy}


\theoremstyle{plain}
%\usepackage{amsmath,amssymb,amsfonts}
%\usepackage{enumerate}
%\usepackage{amscd, amssymb, latexsym, amsmath, amscd}
%\usepackage[all]{xy}
%\usepackage{pb-diagram}

\pagestyle{myheadings}
\newtheorem{theorem}{Theorem}[section]
\newtheorem{thm}[equation]{Theorem}
\newtheorem{prop}[equation]{Proposition}
\newtheorem{cor}[equation]{Corollary}
\newtheorem{con}[equation]{Conjecture}
\newtheorem{lemma}[equation]{Lemma}
\newtheorem{question}[equation]{Question}
\numberwithin{equation}{section}


% The following are only to be used in Math 

\newcommand{\Q}{\mathbb Q}
\newcommand{\Z}{\mathbb Z}
\newcommand{\R}{\mathbb R}
\newcommand{\C}{\mathbb C}
\newcommand{\B}{\mathbb B}
\newcommand{\Sc}{\mathbb S}
\newcommand{\D}{\mathbb D}
\newcommand{\G}{\mathfrak G}
\newcommand{\oO}{\mathcal O}
\newcommand{\F}{\mathbb F}
\newcommand{\h}{\mathbb H}
\newcommand{\hH}{\mathcal H}
\newcommand{\A}{\mathbb A}
\newcommand{\I}{\mathcal A}
\newcommand{\N}{\mathbb N}
\newcommand{\Pp}{\mathbb P}
\newcommand{\E}{\mathcal E}
\newcommand{\s}{\text sym}
%\newcommand{\qed}{\hfill $\Box$ \hfill \\}
\def\M{{\rm M}}

\newcommand{\ad}{{\mathrm{ad}}}
\newcommand{\Aut}{{\mathrm{Aut}}}
\newcommand{\EP}{{\mathrm{EP}}}
\newcommand{\Gal}{{\mathrm{Gal}}}
\newcommand{\Hom}{{\mathrm{Hom}}}
\newcommand{\Ext}{{\mathrm{Ext}}}
\newcommand{\Fr}{{\mathrm{Fr}}}
\newcommand{\rc}{{\mathrm{c}}}
\newcommand{\ri}{{\mathrm{i}}}
\newcommand{\rg}{{\mathrm{gen}}}
\newcommand{\rh}{{\mathrm{h}}}
\newcommand{\rL}{{\mathrm{L}}}
\newcommand{\rM}{{\mathrm{M}}}
\newcommand{\Stab}{{\mathrm{Stab}}}
\newcommand{\sgn}{{\mathrm{sgn}}}
\newcommand{\res}{{\mathrm{Res}}}
\newcommand{\Ind}{{\mathrm{Ind}}}
\newcommand{\uInd}{{\mathrm{^u Ind}}}
\newcommand{\ind}{{\mathrm{ind}}}
\newcommand{\Ann}{{\mathrm{Ann}}}

\newcommand{\rO}{{\mathrm{O}}}
\newcommand{\rU}{{\mathrm{U}}}


\newcommand{\pr}{{\mathrm{pr}}}
\newcommand{\Spin}{{\mathrm{Spin}}}
\newcommand{\Sp}{{\mathrm{Sp}}}
\newcommand{\SL}{{\mathrm{SL}}}
\newcommand{\GL}{{\mathrm{GL}}}
\newcommand{\PGL}{{\mathrm{PGL}}}
\newcommand{\PGSp}{{\mathrm{PGSp}}}
\newcommand{\SU}{{\mathrm{SU}}}
\newcommand{\SO}{{\mathrm{SO}}}
\newcommand{\St}{{\mathrm{St}}}
\newcommand{\Wh}{{\mathrm{Wh}}}



%\usepackage{amsfonts,amssymb,theorem} 
%\usepackage{theorem}
%\usepackage[all]{xy}
\def\e{\epsilon}
\def\gg{{\mathfrak{g}}}
\def\gh{{\mathfrak{h}}}
\def\A{{\mathbb A}}
\def\ws{{\widetilde{\Sp}}}
\def\wpi{{\widetilde{\pi}}}
\def\GG{{\mathbb G}}
\def\NN{{\mathbb N}}
\def\P{{\mathbb P}}
\def\R{{\mathbb R}}
\def\V{{\mathbb V}}
\def\Va{{V_{\alpha}}}
\def\X{{\mathbb X}}
\def\Z{{\mathbb Z}}
\def\T{{\mathbb T}}
\def\gm{{\Gamma}}
\def\vr{{\varphi}}
\def\wx{{\widehat{X}}}
\def\wy{{\widehat{Y}}}
\def\wm{{\widehat{M}}}
\def\gk{{\mathfrak{k}}}
\def\gp{{\mathfrak{p}}}
\def\aa{{\mathfrak{{\huge a}}}}
\def\md{{\medskip}}
\def\ep{{\epsilon}}
\def\bl{{\bullet}}
\def\ms{\:\raisebox{-0.6ex}{$\stackrel{\circ}{B}$}\:}
\def\gs{\:\raisebox{-1.6ex}{$\stackrel{\circ}{A}^c$}\:}
\def\gg{\:\raisebox{-1.6ex}{$\stackrel{\longrightarrow}{\eta \times \eta}$}\:}
\def\ll{\:\raisebox{-0.6ex}{$\stackrel{\lim}{\rightarrow}$}\:}
\def\gk{\:\raisebox{-1.6ex}{$\stackrel{\hookrightarrow}{\rm g}$}\:}
\def\gi{\:\raisebox{-0.6ex}{$\stackrel{\rm dfn}{=}$}\:}
\def\op{\:\raisebox{-0.1ex}{$\stackrel{0}{P}$}\:}
\def\om{\:\raisebox{-0.1ex}{$\stackrel{0}{\Omega}$}\:}
\def\X{\widehat{X}} 
\def\Y{\widehat{Y}} 
%\def\A{{\mathbb A}_{f}}
\def\AA{{\mathbb A}} 
%\def\C{{\bf C}}
\def\T{{\bf T}}
\def\M{{\cal M}}

\def\QQ{{\mathbb Q}}
\def\RR{{\mathbb R}}
\def\PP{{\mathbb P}}
\def\SS{{\mathbb S}}
\def\CC{{\mathbb C}}
\def\ZZ{{\bf Z}}
\def\H{{\cal H}} 
%\def\ll{{\lambda_1}}
\def\L{{\stackrel{m}{\wedge}}}
\def\g{{\bf g}}
\def\b{{\bf b}}
\def\f{{\bf f}}
\def\h{{\bf h}}
\def\k{{\bf k}}
\def\n{{\bf n}}
\def\q{{\bf q}}
\def\u{{\bf u}} 
\def\wp{{W^{\perp}}}
\def\p{{\bf p}}
\def\w{{\bf w}}
\def\CL{{\cal C}}
\def\G{{\mathbb G}} 
\def\wg{{\widehat{G}}}
\def\CT{{\cal T}} 

\setlength{\oddsidemargin}{0.2in}
\setlength{\evensidemargin}{0.2in}
\setlength{\textwidth}{6.3in}


\title {Howe duality for the dual pair $\SL_2(\mathbb R) \times F_{4,1}$ \\ a ping-pong of $K$-types}

\thispagestyle{empty}

\author{Gordan Savin} 
\address{Department of Mathematics, University of Utah, Salt Lake City, UT}
\email{gordan.savin@utah.edu}
\begin{document}


\begin{abstract} 
We prove Howe duality for an exceptional theta correspondence. To that end, 
 we exploit a pair of see-saw identities and relate the $K$-types of corresponding representations.  
\end{abstract} 

\subjclass{22E46, 22E47}
\keywords{minimal representation, theta correspondences} 

\footnote {This work is  supported by the Croatian Science Foundation under the project IP-2022-10-4615 and by a gift No. 946504 from the Simons Foundation.}

\maketitle 



\section{Introduction} 


\vskip 5pt 
Let $\mathbb O$ be the algebra of Cayley octonions over the field of real numbers $\mathbb R$. 
 Let $J$ be the $27$-dimensional space of $3\times 3$ hermitian symmetric matrices with coefficients in $\mathbb O$.  
 Let $N_J : J \rightarrow \mathbb R$ be the cubic form (the norm of $J$), essentially the determinant of $3\times 3$ matrices. 
For every $e\in J$ such that $N_J(e)\neq 0$ there is a structure of exceptional Jordan algebra on $J$ such that $e$ is the identity of $J$. 
Let $G=\Aut(J,e)$ be the group of automorphisms of the resulting Jordan algebra which is the same as the group of linear transformations of $J$ preserving 
$N_J$ and the point $e$.  It is a simple Lie  group of absolute type $F_4$. See \cite{KMRT} for all of this. 
If we pick $e$ to be 
\[ 
\left(\begin{array}{ccc} 
+1 & & \\
 & +1 &  \\
 & & +1 \\
\end{array} 
\right)
\text{ and } 
\left(\begin{array}{ccc} 
-1 & & \\
 & -1 &  \\
 & & +1 \\
\end{array} 
\right). 
\] 
then $G$ is compact and of split rank one, respectively,  for the two choices of $e$ \cite{OV}.   
The Jordan algebra, by way of Koecher-Tits construction \cite{K67}, gives rise to a simply connected group $G(J)$, of the exceptional type $E_7$ and split rank 3 over $\mathbb R$. 
(The same group for both choices of $e$).  The group $G(J)$ comes along with the dual pair (see \cite{KS15}) 
\[ 
\SL_2(\mathbb R) \times G \subset G(J). 
\] 
These dual pairs are completely analogous to $\SL_2(\mathbb R) \times \rO(p,q)$ in $\Sp_{2n}(\mathbb R)$ where $n=p+q$.  Indeed, if we take $J$ to be 
the space of $n\times n$ symmetric matrices with coefficients in $\mathbb R$, then orthogonal groups are stabilizers of generic points in $J$, and 
 $G(J)$ is $\Sp_{2n}(\mathbb R)$.  

\medskip 

The group $G(J)$ has a minimal (holomorphic) representation that appears as a local component of a global representation \cite{Kim}. In \cite{GS}, the 
representation $\Pi$ was restricted to the dual pair $\SL_2(\mathbb R) \times G$, with $G$ compact, and the following decomposiion was obtained: 
\[ 
\Pi= \bigoplus_{n\geq 0} \delta(2n+12) \otimes E_n. 
\] 
Here $\delta(2n+12)$ is the holomorphic representation of the lowest weight $2n+12$ and $E_n$ is the irreducible representation of $G$ of the highest weight 
$n\varpi_4$ where $\varpi_4$ is the fourth fundamental weight for $F_4$. It is the highest weight of the 26-dimensional irreducible representation of $G$ (the complement 
of the line through $e$ in $J$).  

The goal of this short paper is to study the restriction of $\Pi$ to the dual pair with $G$ non-compact. Let $K$ be the maximal compact subgroup of $G$.  
We emphasize that we do not work with continuous representations of 
non compact groups, but with the corresponding $(\mathfrak g, K)$-modules, where $\mathfrak g$ is the complex Lie algebra of $G$. 
 Thus, if $\pi$ is a $(\mathfrak g, K)$-module of finite length, we define 
\[ 
\Theta(\pi)= (\Pi \otimes \pi^{\vee})_{\mathfrak g}. 
\] 
Here $\pi^{\vee}$ is the contragredient of $\pi$, and the subscript $\mathfrak g$ is saying that we are taking co-invaraints with respect to the action of $\mathfrak g$ on 
the tensor product. We can analogously define $\Theta(\sigma)$ for an $(\mathfrak{sl}_2,\SO(2))$-module $\sigma$ of finite length.  Observe that 
$\Theta(\pi)$ and $\Theta(\sigma)$ are  naturally $(\mathfrak{sl}_2,\SO(2))$ and $(\mathfrak g, K)$-modules, respectively.  In this paper we shall show that 
$\Theta(\pi)$ and $\Theta(\sigma)$ are finite length modules, and that they have unique irreducible quotients, if non-zero.  
 The main input is the structure of lifts of types. More precisely, if $\tau$ is a $K$-type, then $\Theta(\tau)$ is also 
 an $(\mathfrak{sl}_2,\SO(2))$-module that we determine explicitly. Similarly, we have a description of the lift for the $\SO(2)$-types.  
This is done in the last section following an idea of Howe \cite{Ho}. 
With these two pieces of information in hand we can play a game of ping pong with types: if $\sigma \otimes \pi$ is a quotient of $\Pi$ and $\tau$ is a type of $\pi$ then, by a 
see-saw identity,  $\sigma$ must have 
an $\SO(2)$-type determined by $\tau$ and vice versa. More details in the next section where main results are obtained.  
A similar strategy (and the name ping-pong) was used  in \cite{GS} to establish Howe duality for exceptional $p$-adic dual pairs. 

\section{Main results} 
The correspondence with compact $G$ establishes a correspondence of infinitesimal characters in the non-compact case. The reader can consult \cite{Li} for more details on this subject. 
Let us write down the correspondence. Using the standard 
realization of the $F_4$ root system, the infinitesimal character of $E_n$ (the representation with the highest weight $n\varpi_4$) is 
\[ 
\frac{1}{2}(2n+11, 5,3,1).  
\] 
On the other hand, the infinitesimal character of $\delta(2n+12)$ is $2n+11$ which we recognize as the first entry above. 
This means, if $\sigma$ has infinitesimal character $x$, then $\Theta(\sigma)$ has infinitesimal character $\frac{1}{2}( x, 5,3,1)$.  More generally, if $\sigma$ is annihilated 
by an ideal in the center of $U(\mathfrak{sl}_2)$ of finite co-dimension, then $\Theta(\pi)$ is also annihilated by an ideal in the center of $U(\mathfrak g)$ of finite co-dimension. 
Hence,  for $\sigma$ of finite length, in order to prove that $\Theta(\sigma)$ has finite length, it suffices to prove that it is admissible. The same goes for $\Theta(\pi)$.   
 
 
\smallskip 


The maximal compact subgroup of $\SL_2(\mathbb R)$ is $\SO(2)$. Its irreducible representations are one-dimensional characters 
parameterized with integers $n$. Let $(n)$ denote the corresponding one-dimensional representation. 
Since the center of $\SL_2(\mathbb R)$ is also the center of 
the simply connected $G(J)$, only even $n=2m$ characters appear in $\Pi$.  

\smallskip 
 
 The maximal compact subgroup of $G$ is denoted by $K$. It is a simple group of 
 type $B_4$. The group $K$ can be picked to be the intersection of $G$ with the compact form of $G$,  
 where the two groups are the stabilizers of  the two choices for $e$, as in the introduction. 
  Let  $\mathfrak g$ be the complex simple Lie algebra of $G$, and 
 $\mathfrak g=\mathfrak k \oplus \mathfrak p$ the  corresponding Cartan decomposition. 
  Here  $\mathfrak p$ is the 16-dimensional spin representation of $K$.  
  For definiteness, assume that the highest weight of $\mathfrak p$ is 
\[ 
\mu=\frac{1}{2}(1,1,1,1). 
\] 
Let $\lambda=(1,0,0,0)$ be the highest weight of the standard 9-dimensional irreducible representation of $K$. 
Let $\tau(m,n)$ be the irreducible representation of $K$ of the highest weight $m\lambda + n\mu$. Only these representations 
of $K$ appear in the restriction of $\Pi$.  This can be seen by applying the branching rule \cite{Lep} to the representations $E_n$.  
 Let $\Theta(\tau(m,n))$ be the lift of $\tau(m,n)$. 
By Proposition \ref{P:1}, we have 
\[  
\Theta(\tau(m,n))\cong U(\mathfrak{sl}_2) \otimes_{U(\mathfrak{so}(2))} \otimes ( 2m+4). 
\] 
A power of this identity is demonstrated by the following lemma:  

\begin{lemma} \label{L:1} 
 Let $\sigma$ be a finite length $(\mathfrak{sl}_2,\SO(2))$-module. Then 
\[ 
\Hom_K(\Theta(\sigma), \tau(m,n)) \cong \Hom_{\mathfrak{sl}_2} (\Theta(\tau(m,n)), \sigma) \cong \Hom_{\SO(2)}( (2m+4) , \sigma). 
\] 
\end{lemma} 
\begin{proof} The first isomorphism is the see-saw identity, obtained by switching the order of taking $\mathfrak{sl}_2$ and $K$ co-invariants. 
The second isomorphism is the Frobenius reciprocity.  
\end{proof} 

Now we have the following consequence, Santa Claus is coming to town: 

\begin{prop} Let $\sigma$ be a finite length $(\mathfrak{sl}_2,\SO(2))$-module. Then $\Theta(\sigma)\neq 0$ if and only if $\sigma$ has a 
type $2m+4$ for some $m\geq 0$.  $\Theta(\sigma)$ has finite length. If $\sigma$ is irreducible, $\Theta(\sigma)$ has multiplicity free 
$K$-types, consisting of all $\tau(m,n)$ such that $2m+4$ is a type of $\sigma$. 
\end{prop} 
\begin{proof} This is all trivial from the lemma, only finite length of $\Theta(\sigma)$ perhaps merits some explanation. 
It is a combination of admissibility (from lemma) and the fact that 
$\Theta(\sigma)$ is annihilated by an ideal in $Z(\mathfrak g)$ of finite co-dimension.  
\end{proof} 

Now we go in the opposite direction. For a character $2m+4$ of $\SO(2)$ consider  $\Theta(2m+4)$. By Proposition \ref{P:2}, 
 it is a quotient of 
\[ 
U(\mathfrak g)\otimes_{U(\mathfrak k)} F_m 
\] 
where $F_m=\mathbb C$ if $m\leq 0$, otherwise  
\[ 
F_m=\tau(0,0) \oplus \tau(1,0) \oplus \cdots \oplus  \tau(m,0).  
\] 

\begin{lemma} Let $\pi$ be a finite length $(\mathfrak g, K)$-module. Then 
\[ 
\Hom_{\SO(2)}(\Theta(\pi), (2m+4) ) \cong \Hom_{\mathfrak g} (\Theta(2m+4), \sigma) \subseteq  \Hom_{K}(F_m, \pi).  
\]
\end{lemma} 
\begin{proof} The first is the see-saw identity, the second is the Frobenius reciprocity.  
\end{proof} 

A consequence of this lemma is that $\Theta(\pi)$ is admissible and hence of finite length. We are now ready to prove the main result. Keep in mind that, without loss 
of generality, we can assume that $\sigma$ has a type $2m+4$ for some $m\geq 0$. (Otherwise the lift is 0.) 


\begin{thm} Let $\sigma$ be an irreducible $(\mathfrak{sl}_2,\SO(2))$-module. Assume that $\sigma$ contains the type $(2m+4)$, for $m\geq 0$, and no smaller types $2n+4$, with $n\geq 0$. 
 (This condition is empty if $m=0$.) Then $\Theta(\sigma)$ has a unique irreducible quotient. It contains the type $\tau(m,0)$ with multiplicity one, and no types $\tau(n,0)$ with $n<m$. 
Conversely, if $\pi$ is an irreducible $(\mathfrak g,K)$-module containing the type $\tau(m,0)$, and 
no smaller such types, then $\Theta(\pi)$ has unique irreducible quotient. It contains the type $(2m+4)$, and no smaller types $2n+4$, with $n\geq 0$. 
\end{thm} 
\begin{proof}  Assume $\pi$ is a quotient of $\Theta(\sigma)$. We do not assume that $\pi$ is irreducible. By Lemma \ref{L:1}, we have the following sequence of inclusions: 
\[ 
\Hom_K(\pi, \tau(m,0)) \subseteq \Hom_K(\Theta(\sigma), \tau(m,0))  \cong \Hom_{\mathfrak{sl}_2} (\Theta(\tau(m,0)), \sigma) \cong \Hom_{\SO(2)} ((2m+4) , \sigma). 
\] 
Observe that we can ran this sequence with any $2n+4$ in place of $2m+4$. If $n<m$, by the assumption, the last space is trivial, hence $\tau(n,0)$ is not a type of 
$\pi$. We shall use this in a moment. 
Since $\sigma$ is a quotient of $\Theta(\pi)$ we have the second sequence of inclusions (note that we are starting with the space of the same dimension as 
$\Hom_{\SO(2)} ((2m+4) , \sigma)$: 
\[ 
\Hom_{\SO(2)}(\sigma, (2m+4)) \subseteq \Hom_{\SO(2)}(\Theta(\pi), (2m+4))  \cong \Hom_{\mathfrak g} (\Theta(2m+4), \pi) \subseteq  \Hom_{K} (F_m , \pi). 
\] 
Since, $\pi$ has no type $\tau(n,0)$ with $n<m$, we ended with  $\Hom_{K} (\tau(m,0) , \pi)$, which has  the same dimension as $\Hom_K(\pi, \tau(m,0))$. 
Thus all inclusions in the two sequences  are isomorphisms, and thus all spaces are one-dimensional, since $\Hom_{\SO(2)} ((2m+4) , \sigma)$ is one-dimensional. 
However, we did not assume that $\pi$ is irreducible. If $\pi=\pi_1\oplus \pi_2$ then, 
\[ 
1+1 = \dim \Hom_{K}(\pi_1,   \tau(m,0)) + \dim\Hom_{K}(\pi_2,  \tau(m,0)) = \dim  \Hom_{K}(\pi, \tau(m,0))=1
\] 
a contradiction. Thus $\Theta(\sigma)$ has a unique irreducible quotient. It contains $\tau(m,0)$, with multiplicity one.  The other direction is proved in the same way... 
\end{proof} 



\section{Computing lifts of types} 

In this section we verify the expressions for $\Theta(\tau(m,n))$ and $\Theta(2m+4)$ used in the proof of the main result. 
The idea comes from \cite{Ho} and uses the following  see-saw diagram in $G(J)$.  
Here $H$ is a simply connected, hermitian symmetric group of absolute type $E_6$.  The centralizer of $H$ is  $\SO(2)$, as the picture indicates. 
The centralizer of $K$ is  $\SL_2\times\SL_2$, where the centralizer of $G$ sits diagonally. 






 \begin{picture}(100,130)(-130,15) 

\put(20,24){$\SO(2)$} 
\put(24,74){$\SL_2$} 
\put(04,124){$\SL_2\times \SL_2$}

\put(74,24){$K$}
\put(75,74){$G$}
\put(75,124){$H$}


\put(37,36){\line(1,2){40}}

\put(42,76){\line(1,0){30}} 

\put(37,116){\line(1,-2){40}}



\end{picture} 

A word of caution here. If we pick a different $\SL_2$  in  $\SL_2\times \SL_2$, the one consisting of all $(g, hgh^{-1})$  where 
$h=\left(\begin{smallmatrix} 1 & 0 \\ 0 & -1 \end{smallmatrix}\right)$, then $G$ and $H$ in the above are replaced by their compact forms.  
In other words, it is important how we identify groups isomorphic to $\SL_2(\mathbb R)$. 

\smallskip 

Let $(e,h,f)$ be an $\mathfrak {sl}_2$-triple such that $\mathbb C \cdot h$ is the Lie algebra of $\SO(2)$.  
For an integer $n>0$, let $\delta(n)$ be the irreducible lowest weight $n$ module. Let $v_n$ be a non-zero lowest weight vector. 
Let $\bar{\delta}(m)$  be the complex conjugate of $\delta(m)$.  It is the irreducible highest weight $-m$ module. 
 Observe that there is a natural map 
\[
U(\mathfrak {sl}_2) \otimes_{U(\mathfrak {so}(2))} (n-m) \rightarrow  \delta(n) \otimes \bar{\delta}(m)
\] 
where $1\in \mathbb C \cong (n-m)$ is mapped to $v_n\otimes v_m$. Since $\delta(n)$ is a free $\mathbb C[e]$-module generated by $v_n$ and 
$\bar{\delta}(m)$ is a free $\mathbb C[f]$-module generated by $v_m$, the above map is easily checked to be an isomorphism.  

\begin{prop} \label{P:1}
Let $\tau(m,n)$ be the irreducible $K$-type of the highest weight $\frac{1}{2}( 2m+n,n,n,n)$.  Then 
\[ 
\Theta(\tau(m,n))\cong U(\mathfrak{sl}_2) \otimes_{U(\mathfrak{so}(2))} \otimes ( 2m+4). 
\] 
\end{prop} 
\begin{proof} Since the centralizer of $K$ is $\SL_2 \times \SL_2$, $\Theta(\tau(m,n))$ is naturally an $\SL_2 \times \SL_2$-module. 
By \cite[Proposition 3.3.3]{Shan} (careful with $\SL_2$'s) we have 
\[ 
\Theta(\tau(m,n))\cong \delta(2m+n+8)  \otimes \bar{\delta}(n+4). 
\] 
In view of the discussion above,  and $(2m+n+8)-(n+4)=2m+4$, the proposition follows. 

\end{proof} 

It remains to discuss $\Theta(2m+4)$. Let $L$ be a maximal compact subgroup of $H$. We can assume that $K\subset L$. Let $\mathfrak h$ and $\mathfrak l$ be the 
complex Lie algebras of $H$ and $L$.  Since $H/L$ is a hermitian symmetric space,  $\mathfrak l$ is a Levi subalgebra such that $[\mathfrak l, \mathfrak l]$ is a simple Lie algebra of type $D_5$. 
We have a Cartan decomposition 
\[ 
\mathfrak h = \bar{\mathfrak u} + \mathfrak l + \mathfrak u
\] 
such that $\mathfrak q= \mathfrak l + \mathfrak u$ is a parabolic subalgebra. If $F$ is a finite-dimensional $\mathfrak l$-modue, we can define a highest weight module 
\[ 
U(\mathfrak h) \otimes_{U(\mathfrak q)} F \cong U(\bar{\mathfrak u}) \otimes F. 
\] 
We now restrict this module to $\mathfrak g=\mathfrak k + \mathfrak p$. Recall that $\mathfrak p$ is 16-dimensional spin-module. On the other hand, $\bar{\mathfrak u}$ and $\mathfrak u$ 
are two 16-dimensional spin modules for $[\mathfrak l, \mathfrak l]$, the simple algebra of type $D_5$. Hence  $\mathfrak p$ must embed diagonally into $\bar{\mathfrak u} + \mathfrak u$. 
Now it is not difficult to check that the natural map 
\[ 
U(\mathfrak g) \otimes_{U(\mathfrak k)} F \rightarrow U(\mathfrak h) \otimes_{U(\mathfrak q)} F
\] 
given by the idenity on $1\otimes F$ is an isomorphism.  We are ready to prove the following: 

\begin{prop} \label{P:2}
Let $m$ be an integer. Then $\Theta(2m+4)$ is a quotient of $U(\mathfrak g)\otimes_{U(\mathfrak k)} F_m$  where 
$F_m=\mathbb C$ if $m\leq 0$, otherwise 
\[ 
F_m=\tau(0,0) \oplus \tau(1,0) \oplus \cdots \oplus \tau(m,0). 
\] 
\end{prop} 
\begin{proof}  $\Theta(2m+4)$ is an $(\mathfrak h, L)$-module, determined in \cite[Section 6]{Pavle}. It is a quotient of the Verma module $U(\mathfrak h) \otimes_{U(\mathfrak q)} F_m$ 
where $F_m$ is a one dimensional representation of $L$ if $m\leq 0$. Otherwise $F_m$, restricted to $[L,L]$, is irreducible with the 
highest weight $(m,0,0,0,0)$.  The restriction of this representation to $K$ is the claimed sum. 
\end{proof} 



%\newpage

\begin{thebibliography}{99}

\bibitem{GS23}
Wee Teck Gan and G. Savin, 
Howe duality and dichotomy for exceptional theta correspondences. 
\textit{Invent. Math.} {\bf 232} (2023), no. 1, 1--78.


\bibitem{GS} B. H. Gross and G. Savin, 
 {\em Motives with Galois group of type $G_2$.}
Compositio Math.  {\bf 114} (1998), 153--217.

\bibitem{Ho} R. Howe, 
 {\em Transcending  classical invariant theory.} 
J. Amer. Math. Soc. {\bf 3}, no 2, (1989), 535--552.  

%\bibitem{KS} E. Karasiewitz and G. Savin,  

\bibitem{Kim} Henry Kim,
 {\em Exceptional Modular form of Weight 4 on an
Exceptional Domain contained in $\mathbb C^{27}$.} 
Revista Matematica Iberoamericana {\bf 9}, No. 1, (1993).

\bibitem{KMRT} M. Knus, A. Merkurjev, M. Rost and J.-P. Tignol,
  {\em The Book of Involutions.}  AMS Colloquium Publications,
  Vol. 44, 1998.
  
  
 \bibitem{KS15} 
T. Kobayashi and G. Savin, 
Global uniqueness of small representations. 
\textit{Math. Z.} {\bf 281} (2015), no. 1-2, 215--239.

\bibitem{K67}
M. Koecher, 
Imbedding of Jordan algebras into Lie algebras. I. 
\textit{Amer. J. Math.} {\bf 89} (1967), 787--816.

\bibitem{Lep} J. Lepowsky, 
 {\em Multiplicity formulas for certain semisimple Lie groups.} 
 Bulletin of the AMS {\bf 77}, no. 4,  (1977) 601--605. 


  
 \bibitem{Li} Jian-Shu Li, {\em The correspondences of infinitesimal
characters for reductive dual pairs in simple Lie groups.} Duke
Math. J. {\bf 97}, no. 2, (1999), 347--377.
  
     
  \bibitem{OV} A. L. Onishchik and E. B. Vinberg (Eds.), 
{\em Lie Groups and Lie Algebras III: Structure of Lie Groups and Lie Algebras}, 
Springer-Verlag, Berlin Heidelberg, 1994.

\bibitem{Pavle} P. Pand\v zi\'c. A. Prli\'c, G. Savin, V. Sou\v cek and V. Tu\v cek, 
{\em On the classification of unitary highest weight modules in the exceptional cases.}  
Journal of Algebra {\bf 684} (2025), 524--562.  

\bibitem{Shan} Yi Shan, 
{\em Exceptional theta correspondence $F_4\times \PGL_2$ for level one automorphic representations.} 
arXiv:2501.19101 
    






\end{thebibliography}

\end{document}



\subsection{Counting representations  of unipotent groups using Bhargava's boxes}
Let $G$ be a split Chevalley group corresponding to a simply laced root system $\Pi$.  Write the highest root $\beta$ as a sum 
$\sum n_{\alpha} \alpha$ of simple roots. Primes dividing the coefficients $n_{\alpha}$ are called bad for  $G$.  Let 
$U$ be the maximal unipotent group in $G$ over the finite field $\mathbb F_p$. A problem of interest to finite group theorists is to 
count irreducible representations of $U$ over the complex field. For example, if $G=\SL_2$, then the number of representations is 
$p= \Phi(p-1)$ where $\Phi(x)=x+1$ is a polynomial with integral and positive coefficients. 

\vskip 5pt 
\underline{Problem}. For every $\Pi$, find a polynomial $\Phi$ with integral and positive coefficients such that the number of irreducible representations 
of $U$ is equal to $\Phi(p-1)$ for all good primes. 

\vskip 5pt 
In particular, although $\Pi$ is simply laced, the representation theory of $U$ may be different for small primes, unless $\Pi$ is of type 
$A_n$. Why is that? One thing that distinguishes the root systems of type $D_n$ and $E_n$ is that the Dynkin diagram has a branch point. 
Let $P$ be the maximal parabolic in $G$ corresponding to the branch point, and let $M$ be its Levi factor. Its derived group is 
$\SL_n \times \SL_m \times \SL_l$. The unipotent radical $U$ has a filtration 
 \[ 
 U \geq U_1 \geq U_2
 \] 
 where $U_1$ is  the unipotent radical of  $P$,  so $U/U_1$ is the  maximal unipotent subgroup of $M$, and 
 \[ 
 U_1/U_2 \cong \mathbb F_p^n \otimes \mathbb F_p^m\otimes \mathbb F_p^l
 \] 
  as $\SL_n \times \SL_m \times \SL_l$-modules. i.e. $U_1/U_2$ is the space of $n\times m\times l$-boxes. In order to count 
  representations of $U$ trivial on $U_2$, by Mackey lemma, one needs to determine orbits of $U/U_1$ on 
$(U_1/U_2)^*\cong U_1/U_2$.  It turns out that bad primes already play a role at this level. Let us illustrate this in the case $D_4$. 
Then $n=m=l=2$ and an element in $U_1/U_2$ is  represented by a cube 

\begin{picture}(100,130)(-130,0) 

\put(20,18){c} 
\put(25,20){\line(1,0){50}}
\put(76,18){d}
\put(22,26){\line(0,1){46}}
\put(78,26){\line(0,1){46}}

\put(40,48){g} 
\put(45,50){\line(1,0){50}}
\put(96,48){h}
\put(42,56){\line(0,1){46}}
\put(98,56){\line(0,1){46}}

\put(25,25){\line(2,3){14}}
\put(82,25){\line(2,3){14}}


\put(20,74){e}
\put(25,76){\line(1,0){50}}
\put(76,74){b}

\put(40,104){a}
\put(45,106){\line(1,0){50}}
\put(96,104){f}

\put(25,81){\line(2,3){14}}
\put(82,81){\line(2,3){14}}


\end{picture}

\noindent 
where vertices are in $\mathbb F_p$.  The action of  an element $(x,y,z)\in \mathbb F_p\oplus \mathbb F_p\oplus \mathbb F_p =U/U_1$ 
is given by the following  three types of ``row-column" operations on cubes: 
\begin{itemize}
\item add $x$ $ \times $ the front face to the rear face of the cube. 
\item add $y$ $\times$  the left face to  the right face of the cube. 
\item add $z$ $\times$  the top face to the bottom face of the cube.
\end{itemize} 
Now it is easy to check that the stabilizer  in $U/U_1$ of the cube with $e=0$ and $abc\neq 0$ consist of triples $(x,y,z)$ such that 
 \[ 
 \left(\begin{array}{ccc} 
 0 & c & b \\
 b & a & 0\\
 c & 0 & a 
 \end{array} 
 \right)
 \left( \begin{array}{c} 
 x \\
 y\\
 z
 \end{array} 
 \right) =
 \left( \begin{array}{c} 
 0 \\
 0\\
 0
 \end{array}
 \right).
 \]
The determinant is $-2abc$ and we discover badness for $p=2$.  In the end, assuming $p\neq 2$, 
one finds that the total number of representations of $U$ is $\Phi(p-1)$ where 
\[ 
\Phi(x)=2x^5 + 15x^4 + 36 x^3 + 34 x^2+12 x + 1. 
\] 
This was computed in \cite{HLM}, however the approach sketched here is simpler, and  can be extended to other groups. 

A goal  of this research is to determine the structure of $U/U_1$-orbits on the space of boxes in the remaining cases, and 
compute the polynomial $\Phi(x)$. 


 
     \vskip 15pt
\begin{thebibliography}{999}

\bibitem[HL]{HLM} F. Himstedt, T. Le and K. Magaard,  
{\em Characters of the Sylow $p$-subgroups of the Chevalley groups of type $D_4(p^n)$.} 
J. Alg {\bf 332} (2011), no 1, 414-427. 




\bibitem[B1]{B1} M. Bhargava, {\em Higher composition laws. I. A new view on Gauss composition}, and quadratic generalizations. Ann. of Math. (2) 159 (2004), no. 1, 217�250.

\bibitem[B2]{B2} M. Bhargava, {\em Higher composition laws. II. On cubic analogues of Gauss composition}, Ann. of Math. (2) 159 (2004), no. 2, 865�886.

\bibitem[B3]{B3} M. Bhargava, {\em Higher composition laws. III. The parametrization of quartic rings}, Ann. of Math. (2) 159 (2004), no. 3, 1329�1360.

\bibitem[KMRT]{KMRT} M. Knus, A. Merkurjev, M. Rost, and J.-P. Tignol,
  {\em The Book of Involutions.}  AMS Colloquium Publications,
  Vol. 44, 1998.

\bibitem[Se]{Se} J.-P. Serre, {\em Galois Cohomology,}  Springer-Verlag, 2002. 



\end{thebibliography}



\end{document}



\section{\bf Dual Pairs in $E_6$}
 We conclude this paper by using the results of this paper to construct certain dual pairs in an exceptional group of type $E_6$.
 \vskip 5pt
 
 \subsection{\bf The group $M_J$.} 
 Let $J$ be a Freudenthal-Jordan $F$-algebra of dimension $9$. Suppose that $J$ corresponds to a pair $(B_K, \tau)$ where $B_K$ is a central simple algebra over an \'etale quadratic $F$-algebra $K$ 
and $\tau$ is an involution of the second kind. Thus, $J =  B_K^{\tau}$ is the subspace of $\tau$-symmetric elements.. 
 \vskip 5pt
 
 The algebra $J$  comes equipped with a cubic norm form $N_J$, and we let
 \[  M_J = \{ g \in \GL(J):   N_J \circ g  = N_J \}.   \]
 In fact,
 \[ M_J = \{ (b, \lambda) \in  B_K^{\times} \times F^{\times}: N_{K/F}(N_{B_K}(g)) \cdot \lambda^3 =1\}  / \{ (t, \lambda) \in F^{\times} \times F^{\times}:  t^2 \cdot\lambda =1\} , \]
 acting on $J = B_K^{\tau}$ via:
 \[  (b, \lambda) \cdot x = \lambda \cdot b x \tau(b). \]
  It contains the automorphism group $\Aut(J)$ as a subgroup, and its Lie algebra is
 \[   Lie(M_J) \cong  \{  b \in B_K:  Tr_{B_K}(b) = 0\}. \]
 \vskip 5pt
 
  \subsection{\bf Groups of type $E_6$.} 
 Now we consider the $F$-vector space
 \[ \mathfrak{g}_J =   \mathfrak{sl}_3 \oplus Lie(M_J) \oplus (F^3 \otimes J) \oplus (F^3 \otimes J)^* \]
 Then $\mathfrak{g}_J$ can be given the structure of a simple Lie algebra of type $E_6$ and we let $G_J = E_6^J$ be the adjoint linear algebraic group with $Lie(G_J)  =\mathfrak{g}_J$.  Then the automorphism group of $G_J$ sits in a short exact sequence of algebraic groups
 \[  \begin{CD}
 1 @>>> G_J @>>> \Aut(G_J) @>>> \Z/2\Z @>>> 1. \end{CD} \]
 This short exact sequence need not split  in general.
 \vskip 5pt
 
 
 \subsection{\bf Dual pair $G_2 \times \Aut(J)$.}
 We can now describe some dual pairs in $G_J$ or rather in $\Aut(G_J)$. It will be easier to do this on the level of Lie algebras. 
 \vskip 5pt
 
 Consider the subspace
 \[  \mathfrak{sl}_3 \oplus F^3 \otimes 1_J \oplus F^3 \otimes 1_J. \]
 This is a Lie sub algebra of type $G_2$. Its centraliser in $\mathfrak{g}_J$ is the Lie subalgebra
 \[  Lie(\Aut(J)) \subset Lie(M_J). \]
 Thus we have a dual pair
 \[  G_2 \times \Aut(J) \subset \Aut(G_J), \]
where we recall that $\Aut(J)$ sits in a short exact sequence
\[  \begin{CD}
1 @>>> \Aut(J)^0 @>>>\Aut(J) @>>> \Z/2\Z @>>> 1. \end{CD} \]
If $J$ is associated to a pair $(B_K, \tau)$, then $\Aut(J)^0 = PGU(B_K, \tau)$ is an adjoint group of type $A_2$. 
\vskip 5pt

\subsection{\bf Dual pair $\Aut(i:E \rightarrow J)  \times G_E$.}

Now we fix an embedding $i: E \longrightarrow J$. Recall that this corresponds to an $E$-twisted composition algebra of $E$-dimension $2$, or equivalently to an $\tilde{M}_E(F)$-orbits o the $E$-twisted Bhargava cube. We have the subgroup
\[ \Aut(i) =  \Aut(i: E \rightarrow J) \subset \Aut(J) \subset M_J, \]
and thus an abelian Lie subalgebra
\[  Lie(\Aut(i)) \subset Lie(M_J). \]
The centraliser of this subalgebra in $Lie(M_J)$ is of the form
\[  Lie(\Aut(i)) \oplus \mathfrak{t}_E \]
for some two dimensional abelian Lie subalgebra $\mathfrak{t}_E$. 

\vskip 5pt

The centraliser of $Lie(\Aut(i))$ in $\mathfrak{g}_J$ 
is of the form 
\[  Lie(\Aut(i)) \oplus \mathfrak{g}_E\] 
with
\[  \mathfrak{g}_E = \mathfrak{sl}_3 \oplus  \mathfrak{t}_E
\oplus F^3 \otimes E \oplus (F^3 \otimes E)^* \]
where $E \hookrightarrow J$ via $i$.
 This Lie algebra is isomorphic to $Lie(G_E)$ (where $G_E$ is the simply connected quasi-split group $Spin_8^E$), and we have
the dual pair 
 \[ \Aut(i: E \rightarrow J) \times  G_E \longrightarrow \Aut(G_J).\]
 Note that this map is not injective.
 
 \vskip 5pt
 
 In particular, we have a see-saw diagram:
 \vskip 10pt
 
 \[
 \xymatrix{
  \Aut(J)  \ar@{-}[dr] \ar@{-}[d] & G_E = {\rm Spin}_8^E   \ar@{-}[d] \\
  \Aut(i: E \rightarrow J) \ar@{-}[ur] & G_2}
\]  
 in $G_J = E_6^J$.
 \vskip 10pt
 
 \noindent In a sequel to this paper, we shall study the theta correspondences of representations and automorphic forms associated to the dual pair $\Aut(i) \times G_E$. 
  
 
